\chapter{The grammar of uncertainty}
\label{ch:counting}
\chaptertag{Sets, permutations, combinations, and inclusion--exclusion}

\begin{question}
Before assigning probabilities, how do we describe precisely what could happen?
\end{question}

A \emph{sample space} is built from possible outcomes; an \emph{event} is a collection of outcomes. Set notation is therefore the grammar of probability.

\section{Set operations}

Let $\Omega$ be a universe of outcomes and $A,B\subseteq\Omega$. The union $A\cup B$ means ``$A$ or $B$ (or both),'' the intersection $A\cap B$ means ``both,'' and the complement $A^c$ means ``not $A$.'' The difference $A\setminus B$ contains outcomes in $A$ but not $B$.

\begin{center}
\begin{tikzpicture}[scale=.95]
  \draw[rounded corners,thick,indigo] (-3,-1.6) rectangle (3,1.8);
  \fill[mint!28] (-.9,.1) circle (1.25);
  \fill[amber!32] (.9,.1) circle (1.25);
  \draw[mint!90!black,thick] (-.9,.1) circle (1.25);
  \draw[amber!85!black,thick] (.9,.1) circle (1.25);
  \node at (-1.25,.25) {$A$};
  \node at (1.25,.25) {$B$};
  \node[above right] at (-3,1.8) {$\Omega$};
  \node[below] at (0,-1.8) {Overlap represents $A\cap B$};
\end{tikzpicture}
\end{center}

De Morgan's laws translate negated statements:
\[
(A\cup B)^c=A^c\cap B^c,
\qquad
(A\cap B)^c=A^c\cup B^c.
\]

\begin{warningbox}
In everyday speech, ``or'' is sometimes exclusive. In probability, $A\cup B$ is inclusive unless explicitly stated otherwise.
\end{warningbox}

\section{Products and sequences}

The Cartesian product $A\times B$ contains ordered pairs $(a,b)$. Repeated trials naturally live in product spaces. Three coin flips have sample space $\{H,T\}^3$, with eight ordered strings. Order matters: $HTT$ and $THT$ are different outcomes.

\section{The multiplication principle}

Start by deciding what one outcome records. An ordered list, an unordered
group, and a list with repetitions are different objects. Before using a
formula, describe when two proposed outcomes count as the same.

The \emph{addition principle} says that counts add over disjoint cases.
For example, a club offers either one of three morning sessions or one of
five afternoon sessions. Choosing exactly one session gives $3+5=8$
possibilities. Choosing one session in \emph{each} period gives $3\cdot5=15$.
The words ``either'' and ``each'' describe different experiments.

If a first choice has $m$ possibilities and, after \emph{each} first choice,
a second choice has $n$ possibilities, there are $mn$ ordered outcomes.
More generally, if the successive numbers of choices are $n_1,\ldots,n_k$,
the number of outcomes is $n_1\cdots n_k$. The available choices may depend
on earlier choices; their \emph{number} at each stage must be the stated one.
If different first choices leave different numbers of continuations, count
each branch and add its contributions instead.

\begin{example}[Codes and officers]
A four-digit code may begin with zero. If digits can repeat, there are
$10^4$ codes. If digits must be distinct, the count is
$10\cdot9\cdot8\cdot7=5040$. Assigning president, secretary, and treasurer
from ten students similarly gives $10\cdot9\cdot8=720$ outcomes when no
student holds two offices.
\end{example}

\section{Permutations and combinations}

There are $n!=n(n-1)\cdots1$ orders of $n$ distinct objects, with $0!=1$.
For $0\le k\le n$, an ordered selection of $k$ distinct objects from $n$
has
\[
  (n)_k=n(n-1)\cdots(n-k+1)=\frac{n!}{(n-k)!}
\]
possibilities. This is sometimes written $P(n,k)$. An unordered selection
of $k$ objects is a \emph{combination}. Every such selection has exactly
$k!$ orderings, so
\[
  \binom nk=\frac{(n)_k}{k!}=\frac{n!}{k!(n-k)!}.
\]
In particular, $\binom n0=\binom nn=1$: there is one empty selection and
one selection containing everything. We use $\binom nk=0$ when $k<0$ or
$k>n$ and $n$ is a nonnegative integer.

\begin{center}
\small
\begin{tabularx}{\textwidth}{@{}XXl@{}}
\toprule
Experiment & What distinguishes outcomes? & Count \\
\midrule
$k$ draws with replacement & Order matters; repetition allowed & $n^k$ \\
$k$ draws without replacement & Order matters; objects distinct & $(n)_k$ \\
A $k$-element subset & Order ignored; no repetition & $\binom nk$ \\
An unordered selection with repetition & Only the number of each type matters & $\binom{n+k-1}{k}$ \\
\bottomrule
\end{tabularx}
\end{center}

\begin{example}[A committee is not a list of officers]
From ten students there are $\binom{10}{3}=120$ three-person committees,
but $720$ assignments to three distinct offices. A committee with one
designated chair has $3\binom{10}{3}=360$ possibilities: the two other
members have no distinct roles. The question about order must be answered
before choosing a formula.
\end{example}

\begin{proposition}[Arrangements with repeated objects]
If $n$ objects consist of $r$ types with multiplicities
$n_1,\ldots,n_r$, where $n_1+\cdots+n_r=n$, the number of distinct
arrangements is $n!/(n_1!\cdots n_r!)$.
\end{proposition}
\begin{proof}
Temporarily distinguish objects of the same type with labels. There are
$n!$ arrangements of the labeled objects. Each arrangement after erasing
the labels comes from exactly $n_1!\cdots n_r!$ labeled arrangements.
Dividing by that constant multiplicity gives the result.
\end{proof}
For example, $\texttt{BANANA}$ has $6!/(3!2!)=60$ distinct arrangements.

\subsection*{Unordered selections with repetition}
Suppose $k$ scoops are chosen from $n\ge1$ flavors, with repetition allowed
and order ignored. An outcome specifies nonnegative counts
$x_1+\cdots+x_n=k$. Arrange $k$ stars and $n-1$ bars in a row: stars
before the first bar count flavor 1, those between the first two bars
count flavor 2, and so on. Empty groups are allowed. Choosing the $k$
star positions among $n+k-1$ positions gives $\binom{n+k-1}{k}$ outcomes.
For two scoops from three flavors $A,B,C$, these are
$AA,AB,AC,BB,BC,CC$, a total of $\binom42=6$.

\begin{warningbox}
Divide a count only when every desired outcome has the same number of
representations. Ordered pairs $AB,BA$ both become $AB$, but $AA$ has
only one ordering. Thus dividing $3^2$ by $2!$ does not count the six
two-scoop selections. The equal-multiplicity argument for combinations
applies to distinct objects.
\end{warningbox}

\section{Binomial coefficients count in two directions}

\begin{proposition}[Symmetry and Pascal's identity]
For $0\le k\le n$,
\[
\binom nk=\binom n{n-k},
\]
and, for $n\ge1$,
\[
\binom nk=\binom{n-1}{k-1}+\binom{n-1}{k}.
\]
\end{proposition}
\begin{proof}
A $k$-element subset determines its $(n-k)$-element complement, proving
symmetry. For Pascal's identity, single out one object. A $k$-element
subset either contains it, leaving $k-1$ choices from the other $n-1$
objects, or omits it, leaving $k$ choices from those objects. These two
cases are disjoint and exhaustive.
\end{proof}

\begin{theorem}[Binomial theorem]
For a nonnegative integer $n$,
\[
  (x+y)^n=\sum_{k=0}^n\binom nk x^{n-k}y^k.
\]
\end{theorem}
\begin{proof}
Expand the product of $n$ copies of $(x+y)$. To produce $x^{n-k}y^k$,
choose the $k$ factors that supply $y$; all the others supply $x$.
There are $\binom nk$ choices.
\end{proof}

\begin{example}[Coefficients carry both counts and weights]
In $(2-x)^4$, choosing $k$ copies of $-x$ and $4-k$ copies of $2$ gives
\[
 (2-x)^4=16-32x+24x^2-8x^3+x^4.
\]
For instance, the $x^2$ coefficient is $\binom42 2^2(-1)^2=24$.
The binomial coefficient counts the positions; the remaining factors
give the contribution from each choice.
\end{example}

The same distinction explains a probability formula. In $n$ independent
trials with success probability $p$, a \emph{specified} string with $k$
successes has probability $p^k(1-p)^{n-k}$. There are $\binom nk$ such
strings, so the probability of exactly $k$ successes is
$\binom nk p^k(1-p)^{n-k}$. Their sum is
$(p+(1-p))^n=1$. Chapter 5 develops this binomial distribution; for
$p=0$ or $1$, every trial has its certain outcome.

Taking $x=y=1$ gives $\sum_{k=0}^n\binom nk=2^n$. This also counts all
subsets in two ways: by their sizes, or by making an include/exclude
choice for each object. Taking $x=1,y=-1$ gives, for $n\ge1$,
$\sum_{k=0}^n(-1)^k\binom nk=0$. This alternating sum is the engine
behind the general inclusion--exclusion formula.

\section{Inclusion--exclusion for finite sets}

\begin{proposition}[Two and three sets]
For finite sets $A,B,C$,
\begin{align*}
|A\cup B|&=|A|+|B|-|A\cap B|,\\
|A\cup B\cup C|&=|A|+|B|+|C|\\
&\quad-|A\cap B|-|A\cap C|-|B\cap C|
  +|A\cap B\cap C|.
\end{align*}
\end{proposition}

For two sets, subtraction repairs the double count of the overlap. With
three sets, an element in all three is added three times and subtracted
three times, so it must be added back once. An element in exactly two
sets contributes $2-1=1$; an element in exactly one contributes $1$.
An element outside the union contributes zero. This proves the formulas.

\begin{warningbox}
$A\cap B$ includes the triple intersection $A\cap B\cap C$.
It does not mean ``in $A$ and $B$ but outside $C$.'' Confusing these
two regions produces incorrect inclusion--exclusion calculations.
\end{warningbox}

\begin{example}[Three activities]
Among 100 students, 39 attend activity $A$, 29 attend $B$, and 25 attend
$C$. The pairwise intersection counts are 13, 11, and 9, respectively,
and 5 students attend all three. Thus the union has
$39+29+25-13-11-9+5=65$ students; 35 attend none.
The exactly-two count is $(13-5)+(11-5)+(9-5)=18$.
The exactly-one count is $65-18-5=42$.
\end{example}

\begin{theorem}[Finite inclusion--exclusion]
For finite sets $A_1,\ldots,A_r$,
\[
\left|\bigcup_{i=1}^r A_i\right|
=\sum_{\varnothing\ne J\subseteq\{1,\ldots,r\}}
 (-1)^{|J|+1}\left|\bigcap_{j\in J}A_j\right|.
\]
Equivalently: add the single-set sizes, subtract the pairwise
intersection sizes, add the triple-intersection sizes, and continue
with alternating signs through the intersection of all $r$ sets.
\end{theorem}
\begin{proof}
An element belonging to exactly $t\ge1$ sets occurs in
$\binom tk$ of the $k$-fold intersections. Its total contribution is
\[
\sum_{k=1}^t(-1)^{k+1}\binom tk=1-(1-1)^t=1.
\]
An element in no set contributes zero. Thus each element of the union
is counted exactly once.
\end{proof}

\begin{example}[Using all three symbols]
There are $3^4$ strings of length four over $\{A,B,C\}$.
To count those using every symbol, subtract strings missing at least
one symbol. Each missing-symbol event has $2^4$ strings; each pair of
missing-symbol events has one string; no length-four string misses
all three symbols. Therefore the required count is
\[
3^4-3\cdot2^4+3\cdot1^4=36.
\]
These are counts; turning them into probabilities requires a specified
probability model.
\end{example}

\begin{example}
A five-card poker hand is an unordered choice from 52 cards, so there are $\binom{52}{5}$ hands. Exactly two aces means choosing two of four aces and three of the remaining 48 cards:
\[
\binom42\binom{48}{3}.
\]
\end{example}

\begin{labbox}
Use the website's counting lab to compare committees, ordered selections,
and repeated draws. Then use the inclusion--exclusion lab to vary the eight
disjoint regions of three sets, including the region outside the union.
Recover the marginal and intersection totals from those regions. With
all seven marginal/intersection totals specified, the seven interior
region counts are determined; knowing the universe size also determines
the outside region. If totals are omitted, the remaining data need not
determine a unique diagram.
\end{labbox}

\begin{exercisebox}
\begin{enumerate}
  \item Write $(A\cap B^c)\cup(A^c\cap B)$ in words. \exlevel{Fluency}
  \item Verify both De Morgan laws using membership tables. \exlevel{Reasoning}
  \item How many eight-character strings contain exactly three letters and five digits? Use 26 letters and 10 digits, allow repetition, and allow a leading zero. \exlevel{Fluency}
  \item Count five-card hands containing at least one ace. \exlevel{Reasoning}
  \item Prove $\binom nk=\binom n{n-k}$ by a counting argument. \exlevel{Reasoning}
  \item Among integers from 1 to 1,000, how many are divisible by 2 or 3 or 5? \exlevel{Reasoning}
  \item Design two different sampling procedures that produce the same set of possible outcomes but not the same probabilities. \exlevel{Exploration}
  \item From nine students, count a four-person committee, a list of four distinct officers, and a four-person committee with one chair. \exlevel{Fluency}
  \item Count the distinct arrangements of the letters in $\texttt{LEVEL}$. Explain the divisor in your formula. \exlevel{Fluency}
  \item Prove Pascal's identity by deciding whether a distinguished object is selected. \exlevel{Reasoning}
  \item Find the coefficient of $x^2$ in $(1+3x)^5$, and explain its binomial coefficient. \exlevel{Fluency}
  \item Count length-five strings over $\{A,B,C\}$ that use all three symbols. \exlevel{Reasoning}
  \item In five independent trials with success probability $1/3$, find
  the probability of exactly three successes. Distinguish one specified
  string from all strings with three successes. \exlevel{Fluency}
  \item Choose four scoops from three flavors, allowing repetition.
  Count ordered lists and unordered selections. Explain why dividing
  the first count by $4!$ is invalid. \exlevel{Reasoning}
  \item Count the distinct arrangements of $\texttt{MISSISSIPPI}$.
  Identify every repeated-letter factor. \exlevel{Fluency}
  \item Prove $\sum_{k=0}^n k\binom nk=n2^{n-1}$ for $n\ge1$ by
  counting subsets with one distinguished member in two ways.
  Appendix~\ref{app:calculus} also obtains this identity by differentiation.
  \exlevel{Reasoning}
\end{enumerate}
\end{exercisebox}
